\documentclass[11pt]{article}
\usepackage[margin=1in]{geometry}
\usepackage[T1]{fontenc}
\usepackage[utf8]{inputenc}
\usepackage{hyperref}
\title{orchestrationTools.wl Module Documentation}
\author{Manus AI}
\date{}
\begin{document}
\maketitle

\section{High-level explanation}

\texttt{orchestrationTools.wl} is the high-level workflow glue for the active MFGraphs package. It does not introduce new mathematical data structures. Instead, it connects the existing scenario, unknown, system, and solver layers into convenient top-level calls.

This file is important because it presents the cleanest end-to-end public story of the current package. A user with a scenario object typically wants a solution, not a manual reconstruction of all intermediate stages. This module is what turns the staged architecture into a practical workflow.

\section{Low-level explanation of functions and functionality}

\subsection*{\texttt{solveScenario}}

\texttt{solveScenario[s, solver]} is the main orchestration entrypoint. On a single scenario, it builds the symbolic unknown bundle with \texttt{makeSymbolicUnknowns}, builds the structural system with \texttt{makeSystem}, and then applies the selected solver. If no solver is supplied, it defaults to \texttt{dnfReduceSystem}.

The function also has a list overload: \texttt{solveScenario[\{s1,s2,\ldots\}, solver]} maps the same workflow across multiple scenarios and returns a corresponding list of solutions.

Typical usage is:

\begin{verbatim}
sol = solveScenario[s];
sol2 = solveScenario[s, activeSetReduceSystem];
sols = solveScenario[{s1, s2}];
\end{verbatim}

\subsection*{\texttt{SolveMFG}}

\texttt{SolveMFG[assoc]} is a backward-compatibility wrapper for legacy raw-association input. It first calls \texttt{makeScenario} to convert the raw input into the current typed scenario representation. If conversion fails, it emits a message and returns the failure. Otherwise, it delegates to \texttt{solveScenario}.

Its main role is transitional. New code should generally prefer explicit \texttt{makeScenario} plus \texttt{solveScenario}, but \texttt{SolveMFG} preserves a familiar entrypoint for older calling styles.

\section{Usage guidance}

For most end users, this file contains the friendliest solve API in the repository. If the caller already has a scenario, \texttt{solveScenario} is usually the right choice. If the caller is migrating older code based on raw associations, \texttt{SolveMFG} provides a compatibility bridge.

\section*{References}

\begin{enumerate}
\item Source file: \texttt{MFGraphs/orchestrationTools.wl}
\item Upstream dependencies: \texttt{scenarioTools.wl}, \texttt{unknownsTools.wl}, \texttt{systemTools.wl}, \texttt{solversTools.wl}
\end{enumerate}

\end{document}
